% Pursuing Stacks (cotes 134-2 to 134-8): archive cote 134-4, site page 41 (PDF page 42), typescript page 129.
% typescript. Draft compared: ps3.tex:2044-2124.
% AI-transcribed by Claude Opus 5.5 (claude-opus-5-5, adaptive thinking, effort high), Message Batches API, batch ps2 (msgbatch_01NVwr1Yuo8qJf5rykeGBXPK), 2026-11-12; image review recorded in records/CHECK_DECISIONS_134_3_4.tsv.
\scanpage{134-4/41}
\mspage{129}

% Hand-drawn arrows are marked blue; typed objects remain plain.
\diagram{\begin{tikzcd}[column sep=large]
Q(M) \arrow[color=blue, r, bend left=15, "\hand{r}"] \arrow[color=blue, dd, bend right=50, "\hand{}"'] & L(M) \arrow[color=blue, l, bend left=15, "\hand{s}"] \\
\hand{Q_0=}\mathrm{Hom}_1(M) \arrow[color=blue, u, hook] \arrow[color=blue, r, leftrightarrow, "\hand{\sim}"] & \mathrm{Hom}_2(M)\hand{=L_0} \arrow[color=blue, u, hook] \\
\mathrm{Hom}_3(M) \arrow[color=blue, u, hook] & \mathrm{Hom}_4(M) \arrow[color=blue, u, hook] \arrow[color=blue, l, hook']
\end{tikzcd}}
\hand{\}} (``homotopy structures'' on M,
defined in terms of homotopy
\deleted{\unsure{notions}} relations\hand{,} \over{or}{or}\hand{\ensuremath{\circ}}\diagram{\begin{tikzpicture}[baseline=-.3ex,x=1em,y=1ex]\draw[color=blue,line width=.5pt] (0,-.6)--(0,2.8);\draw[color=blue,line width=.5pt] (-.13,-.1)--(.14,.05);\end{tikzpicture}} homotopism stru
% NOTE: the right edge of the scan cuts the end of the line "homotopism stru..."
tures )
\margin{\deleted{\unsure{déployer}} déployer le diagramme, trop tassé dans les deux dimensions}

\begin{tabular}{@{}l@{\quad}c@{\quad}l@{}}\shortstack[l]{(weak homotopy\\interval struc-\\tures)} & \hand{\vrule width .4pt height 4ex depth 1.2ex} & \shortstack[l]{(contractibility\\structures)}\end{tabular}
,

which in terms of the preceding reflections, and yesterday's, is self-%
explanatory. The vertical\xout{\unsure{l}} arrow from Q(M) to $\mathrm{Hom}_3(M)$ is the canonical
retraction - in terms of the latter, and it's composition with
$s\colon L(M)\hand{\to} Q(M)$, there \uline{are} ways after all to associate to any $R$ in $Q(M)$
or $W$ in $L(M)$ a homotopy structure on $M$, provided only $M$ admits finite
products, \inserted{by} using weak homotopy intervals. If $h$ is the homotopy structure
thus defined, we get a priori a functor
\[ M_h \hand{\to} M_R = M/R \qquad \text{resp.} \quad M_h \hand{\to} M_W = W^{-1}M \quad , \]
\margin{$\times$} provided we assume $\over[typed]{R}{R}$ resp. $W$ satisfy the ``accessory assumption'', namely
that the corresponding functor $M\hand{\to} M'$ (where $M'$ is either $M_R$ or $M_W$)
commute with finite products.

T\over[typed]{h}{H}e main fact to remember from this whole discussion, it seems
to me, is that there are not really \uline{four}, but \uline{only} \uline{three} essentially
distinct types of structure \inserted{(among yesterday's)} we may \xout{\unsure{puton}} consider upon M as \hand{``}homotopy-%
flavored\hand{''} structures, namely
\[
\begin{array}{ccccc}
\text{homotopy structures} & \hand{\Leftarrow} & \text{weak homotopy interval} & \hand{\Leftarrow} & \text{contractibility}\\
 & & \text{structures} & & \text{structures}\\
\mathrm{Hom}_1(M) \hand{\leftrightarrow} \mathrm{Hom}_2(M) & & \mathrm{Hom}_3(M) & & \mathrm{Hom}_4(M).
\end{array}
\]
It would seem at present that the homotopy structures \over[typed]{t}{w}hich naturally
come up in our present ``modelizer story'' are all of the strictest type,
and even describable in terms of just one \uline{generating contractible weak
homotopy interval} (I would like to drop the qualification ``weak'' definit\hand{e-}%
ly, when a contractibility assumption comes in !), and even a \uline{generating
contractor}, with commutative idempotent composition law.

\margin{41}
